% post_recovery_plan Phase 4 (P4-5) example.
%
% Demonstrates the nrpylatex math-grammar integration: % mechanics
% directives describe the FE problem (mesh, formulation, material,
% boundary), and a $...$ math block carries an indexed-tensor
% expression that the new frontend.math_parser routes through
% mechdsl.symbolic.bridge into a SymbolicNode map.
%
% Compile path:
%   from pathlib import Path
%   from mechdsl.frontend import parse_with_math
%   src = Path("examples/svk_latex_math.tex").read_text()
%   context = parse_with_math(src)
%   # context["math"]["tensors"] now carries SymbolicNode entries.
%
% Note (deferral): the closed-form SVK first Piola-Kirchhoff stress
% requires \det F and \log J intrinsics that nrpylatex 1.4.0 does
% not register. The $...$ block below uses the rank-2 copy
% A^{i I} = F^{i I} as the import-chain surrogate; the bridge
% nonetheless preserves the spatial/material index distinction
% (07-CONVENTIONS) so downstream layers see a true two-point tensor.

% mechanics dim 3
% mechanics cell hex8
% mechanics formulation total_lagrangian
% mechanics material svk --E 200e3 --nu 0.3
% mechanics boundary Gamma_u --type dirichlet --value 0
% mechanics boundary Gamma_t --type neumann --traction "t_bar"

% nrpylatex declarations for the math block:
% declare FUU --dim 3
% declare AUU --dim 3

% First-Piola-flavoured surrogate (see header note):
$A^{i I} = F^{i I}$
